\begin{abstract}
Alignment techniques such as supervised fine-tuning and reinforcement learning from human feedback are typically validated in single-turn or short multi-turn settings, yet real-world deployment involves extended conversations spanning tens or hundreds of turns. We investigate whether alignment effects degrade over conversation depth---specifically, whether aligned language models regress toward their base (pre-training) prior as turn count increases. We conduct two complementary experiments: (1)~token-level distribution analysis comparing \qwenbase and \qweninst across conversations of 0--12 turns, and (2)~behavioral alignment probes measuring safety refusal, formatting compliance, and instruction retention in \gptmini and \gptnano across 0--24 turns. We find that superficial alignment behaviors---formatting constraints and style instructions---degrade significantly after 10--15 turns (uppercase compliance drops from 100\% to 25\%, $p<0.001$, Cohen's $d=2.45$), while safety alignment remains robust through 24 turns. Token-level analysis reveals a two-phase pattern: alignment signal \emph{strengthens} in early turns as chat formatting activates, then gradually weakens from turns 2--12. Multi-turn conversation structure itself produces systematically higher instruct--base divergence compared to concatenated equivalents ($\Delta \approx 3.1$ nats, $p<0.001$), suggesting that turn structure interacts with alignment mechanisms in unexpected ways. These findings demonstrate that alignment is not monolithic: different alignment behaviors occupy different depths and exhibit different durability profiles under extended interaction.
\end{abstract}
